%=====================================================================
\section{Three Variational Theorems on $\partial(\Delta^5)$}
\label{ch:variational_theorems}
%=====================================================================

The boundary of the 5-simplex $\partial(\Delta^5)$ admits a unique
$(3,3)$ vertex partition into A-type (spatial) and B-type (temporal)
vertices (Proposition~\ref{prop:33_partition}).  Each of the three
nuclear simplices $\sigma_3, \sigma_4, \sigma_5$ contains the full
A-triple $\{A_1, A_2, A_3\}$ and a B-pair chosen from
$\{B_1, B_2, B_3\}$.  Because $\dim(\CC^2) = 2$, the third B-vertex
satisfies (Proposition~\ref{prop:B3})
\begin{equation}\label{eq:B3_parametrize}
B_3 = \cos\varphi\; B_1 + \sin\varphi\; B_2,
\qquad \varphi \in (0, \pi/2),
\end{equation}
where $B_1, B_2$ are orthonormal: $\langle B_1 | B_2 \rangle = 0$.
We work throughout with the block-diagonal Gram matrix
\begin{equation}\label{eq:block_gram}
G = \begin{pmatrix} G_A & 0 \\ 0 & G_B \end{pmatrix},
\qquad
G_A = \langle A_i | A_j \rangle_{3\times 3},
\quad
G_B = \langle B_i | B_j \rangle_{3\times 3},
\end{equation}
encoding the orthogonality $\langle A_i | B_j \rangle = 0$ for all
$i, j$.  This is the $(3,2)$ block structure restricted to each
nuclear simplex.

This chapter proves three theorems that follow rigorously from this
structure.  Theorems~1 and~2 hold for \emph{all} values of $\varphi$;
Theorem~3 determines $\varphi$ via the variational principle.
Together they supply the three constants entering the closed propagator
(Chapter~\ref{ch:masses}): the alternating sign, the impedance, and
the lattice speed of light.

%---------------------------------------------------------------------
\subsection{Theorem 1: The AAA deficit angle}
\label{sec:thm_AAA}
%---------------------------------------------------------------------

\begin{theorem}[AAA deficit angle]\label{thm:delta_AAA}
Let $\partial(\Delta^5)$ carry the $(3,3)$ partition with
block-diagonal Gram matrix~\eqref{eq:block_gram} and
$B_3 = \cos\varphi\; B_1 + \sin\varphi\; B_2$.  Then the deficit
angle at the unique AAA hinge $\{A_1, A_2, A_3\}$ is
\begin{equation}
\boxed{\delta_{AAA} = \pi}
\end{equation}
independent of $\varphi$.
\end{theorem}

\begin{proof}
The AAA hinge belongs to the three nuclear simplices
$\sigma_3, \sigma_4, \sigma_5$, obtained by omitting $B_1, B_2, B_3$
respectively.  Each nuclear simplex has five vertices: the A-triple
plus a B-pair.  We must compute the dihedral angle of each simplex at
the AAA hinge.

Write $\alpha = \cos\varphi$, $\beta = \sin\varphi$, so that
$\alpha^2 + \beta^2 = 1$.

\medskip\noindent\textbf{Step 1: Factorization.}
Because the Gram matrix is block-diagonal, the $5\times 5$ Gram
determinant of each simplex factorizes:
\[
\det(G_\sigma) = \det(G_A) \cdot \det(G_{B\text{-pair}}).
\]
The cofactors appearing in the dihedral angle formula likewise
factorize.  In the Regge calculus expression for the dihedral angle at
a codimension-2 hinge, the A-sector Gram matrix $G_A$ cancels
identically from the numerator and denominator, leaving the angle
determined entirely by the B-sector.

\medskip\noindent\textbf{Step 2: B-sector overlaps.}
The three B-pairs and their overlaps are:
\begin{align}
\sigma_3 &= \{A_1, A_2, A_3, B_2, B_3\}:
  &\langle B_2 | B_3 \rangle &= \beta, \label{eq:overlap_23}\\
\sigma_4 &= \{A_1, A_2, A_3, B_1, B_3\}:
  &\langle B_1 | B_3 \rangle &= \alpha, \label{eq:overlap_13}\\
\sigma_5 &= \{A_1, A_2, A_3, B_1, B_2\}:
  &\langle B_1 | B_2 \rangle &= 0. \label{eq:overlap_12}
\end{align}
After the factorization of Step~1, the dihedral angle of $\sigma_k$
at the AAA hinge reduces to $\theta_k = \arccos\bigl(\text{Re}\,
\langle B_j | B_l \rangle\bigr)$, where $\{B_j, B_l\}$ is the B-pair
of $\sigma_k$.  Hence:
\begin{align}
\theta_3 &= \arccos(\beta) = \arccos(\sin\varphi), \\
\theta_4 &= \arccos(\alpha) = \arccos(\cos\varphi) = \varphi, \\
\theta_5 &= \arccos(0) = \frac{\pi}{2}.
\end{align}

\medskip\noindent\textbf{Step 3: Summation.}
Using the identity $\arccos(\sin\varphi) = \pi/2 - \varphi$ for
$\varphi \in [0, \pi/2]$:
\begin{equation}
\sum_{k=3}^{5} \theta_k
= \underbrace{\left(\frac{\pi}{2} - \varphi\right)}_{\theta_3}
+ \underbrace{\varphi}_{\theta_4}
+ \underbrace{\frac{\pi}{2}}_{\theta_5}
= \pi.
\end{equation}
The deficit angle is
\begin{equation}
\delta_{AAA} = 2\pi - \sum\theta_k = 2\pi - \pi = \pi. \qedhere
\end{equation}
\end{proof}

\begin{remark}[Universality]
The parameter $\varphi$ cancels algebraically at the level of the
arccos identity.  No specific value of $\varphi$ is needed; the result
holds for all unit-norm decompositions of $B_3$.
\end{remark}

\begin{remark}[Numerical verification]
EXP-047b confirms $\delta_{AAA} = 180.0000^\circ$ at
$\varphi = 10^\circ, 30^\circ, 45^\circ, 60^\circ, 80^\circ$.
\end{remark}

\begin{corollary}[Alternating sign in the propagator]
\label{cor:alternating}
The Regge holonomy around the AAA hinge is
\begin{equation}
e^{i\delta_{AAA}} = e^{i\pi} = -1.
\end{equation}
A fermion pattern encircling this hinge picks up a factor $-1$ per
loop.  The self-energy series therefore alternates:
\begin{equation}
\Sigma = x - x^2 + x^3 - \cdots = \frac{x}{1+x},
\end{equation}
yielding the closed propagator
\begin{equation}
P(x) = 1 + \Sigma = \frac{1 + 2x}{1 + x},
\end{equation}
where $x = \agut \times f_{\text{sector}}$ (Chapter~\ref{ch:masses}).
The convergence is guaranteed: $|x| < 1$ for all physical sectors.
\end{corollary}

%---------------------------------------------------------------------
\subsection{Theorem 2: The mean ABB determinant}
\label{sec:thm_ABB}
%---------------------------------------------------------------------

\begin{theorem}[Mean ABB hinge determinant]\label{thm:det_ABB}
Under the same conditions as Theorem~\ref{thm:delta_AAA},
the average Gram determinant over the 9 ABB-type hinges of
$\partial(\Delta^5)$ is
\begin{equation}
\boxed{\bigl\langle \det(G_h) \bigr\rangle_{ABB} = \frac{n_B}{n_A}
= \frac{2}{3}}
\end{equation}
independent of $\varphi$.
\end{theorem}

\begin{proof}
An ABB hinge has the form $\{A_i, B_j, B_k\}$ with $1 \leq i \leq 3$
and $\{j,k\} \subset \{1,2,3\}$, $j < k$.  The total count is
$\binom{3}{1}\binom{3}{2} = 9$, decomposing as $3$~A-vertices
$\times$ $3$~B-pairs.

\medskip\noindent\textbf{Step 1: Block factorization of the determinant.}
Because $\langle A_i | B_j \rangle = 0$, the $3 \times 3$ Gram matrix
of each ABB hinge is block-diagonal:
\begin{equation}
G_h = \begin{pmatrix} 1 & 0 & 0 \\ 0 & 1 & \langle B_j | B_k\rangle \\
0 & \overline{\langle B_j | B_k\rangle} & 1 \end{pmatrix}.
\end{equation}
The $1\times 1$ A-block contributes $\det = 1$.  Hence
\begin{equation}\label{eq:det_ABB_factor}
\det(G_h) = 1 \cdot \det\!\begin{pmatrix} 1 & \langle B_j|B_k\rangle \\
\overline{\langle B_j|B_k\rangle} & 1 \end{pmatrix}
= 1 - |\langle B_j | B_k \rangle|^2.
\end{equation}

\medskip\noindent\textbf{Step 2: B-pair overlaps.}
The three B-pairs yield:
\begin{center}
\begin{tabular}{lccc}
\toprule
B-pair & $|\langle B_j | B_k \rangle|^2$ & $\det(G_h)$ & Multiplicity \\
\midrule
$\{B_1, B_2\}$ & $0$ & $1$ & $\times 3$ \\
$\{B_1, B_3\}$ & $\cos^2\varphi$ & $\sin^2\varphi$ & $\times 3$ \\
$\{B_2, B_3\}$ & $\sin^2\varphi$ & $\cos^2\varphi$ & $\times 3$ \\
\bottomrule
\end{tabular}
\end{center}
The multiplicity of 3 for each B-pair arises from the three choices
of A-vertex $A_i$.

\medskip\noindent\textbf{Step 3: Average.}
\begin{align}
\bigl\langle \det(G_h) \bigr\rangle_{ABB}
&= \frac{1}{9}\Bigl[3 \cdot 1 + 3\sin^2\varphi + 3\cos^2\varphi\Bigr]
\notag\\
&= \frac{1}{9}\Bigl[3 + 3\bigl(\sin^2\varphi + \cos^2\varphi\bigr)\Bigr]
\notag\\
&= \frac{1}{9}\bigl[3 + 3\bigr]
= \frac{6}{9} = \frac{2}{3}. \qedhere
\end{align}
\end{proof}

\begin{remark}[Universality]
The Pythagorean identity $\sin^2\varphi + \cos^2\varphi = 1$ is the
sole ingredient beyond the block structure.  Like Theorem~1, the
result is independent of the mixing angle $\varphi$.
\end{remark}

\begin{remark}[Numerical verification]
EXP-047b confirms $\langle\det(G_h)\rangle_{ABB} = 0.666667$ at all
tested values of $\varphi$.
\end{remark}

\begin{corollary}[Propagation impedance]\label{cor:impedance}
The propagation impedance of an ABB hinge is defined as the inverse of
the mean hinge determinant:
\begin{equation}
\rho \;=\; \frac{1}{\langle\det(G_h)\rangle_{ABB}}
\;=\; \frac{n_A}{n_B} \;=\; \frac{3}{2}.
\end{equation}
This quantity enters the lepton mass ratio
(Theorem~\ref{thm:lepton_mass}):
\begin{equation}
\frac{m_\mu}{m_e} = \frac{\rho}{\alpha}
\bigl(1 + \text{corrections}\bigr)
= \frac{3}{2\alpha}
\bigl(1 + \text{corrections}\bigr)
= 206.80,
\end{equation}
in agreement with the observed value $206.768$ to $0.02\%$.
\end{corollary}

%---------------------------------------------------------------------
\subsection{Theorem 3: The variational overlap}
\label{sec:thm_overlap}
%---------------------------------------------------------------------

Theorems~1 and~2 hold for all $\varphi$.  The variational principle
$\delta S / \delta\psi = 0$ on the Regge action selects a definite
value.

\begin{theorem}[Variational overlap]\label{thm:overlap}
Let $S(\varphi)$ be the Regge action of $\partial(\Delta^5)$ with the
$(3,3)$ block-diagonal Gram matrix and
$B_3 = \cos\varphi\; B_1 + \sin\varphi\; B_2$.  Then
$\varphi = \pi/4$ is a stationary point of $S$, and
\begin{equation}
\boxed{|\langle B_1 | B_3 \rangle|^2 = \cos^2\!\left(\frac{\pi}{4}\right) = \frac{1}{2}.}
\end{equation}
\end{theorem}

\begin{proof}
The proof proceeds in two steps: a symmetry argument establishes
stationarity, and a numerical computation confirms maximality.

\medskip\noindent\textbf{Step 1: Exchange symmetry.}
Consider the exchange $B_1 \leftrightarrow B_2$.  In
$\partial(\Delta^5)$, the vertices $B_1$ and $B_2$ enter
symmetrically: they are both orthonormal B-type vectors, and no
structure in the Gram matrix~\eqref{eq:block_gram} distinguishes
them (the choice of which is ``$B_1$'' and which is ``$B_2$'' is a
gauge convention).

Under $B_1 \leftrightarrow B_2$, the parametrization
$B_3 = \cos\varphi\; B_1 + \sin\varphi\; B_2$ transforms as
$\varphi \mapsto \pi/2 - \varphi$.  This exchanges
\begin{equation}
\sigma_3 = \{A_1,A_2,A_3,B_2,B_3\}
\;\longleftrightarrow\;
\sigma_4 = \{A_1,A_2,A_3,B_1,B_3\},
\end{equation}
while $\sigma_5 = \{A_1,A_2,A_3,B_1,B_2\}$ is invariant.  Since the
Regge action is a sum over all simplices, it is invariant under
permutations of the summands:
\begin{equation}\label{eq:S_symmetry}
S(\varphi) = S(\pi/2 - \varphi).
\end{equation}

\medskip\noindent\textbf{Step 2: Stationarity at $\varphi = \pi/4$.}
Differentiating~\eqref{eq:S_symmetry} with respect to $\varphi$:
\begin{equation}
\frac{dS}{d\varphi} = -\frac{dS}{d\varphi}\bigg|_{\pi/2-\varphi}.
\end{equation}
Evaluating at $\varphi = \pi/4$ (the fixed point of
$\varphi \mapsto \pi/2 - \varphi$):
\begin{equation}
\frac{dS}{d\varphi}\bigg|_{\pi/4}
= -\frac{dS}{d\varphi}\bigg|_{\pi/4},
\end{equation}
which implies
\begin{equation}
\frac{dS}{d\varphi}\bigg|_{\pi/4} = 0.
\end{equation}
Hence $\varphi = \pi/4$ is a stationary point of the Regge action.

\medskip\noindent\textbf{Step 3: Maximum (numerical).}
Numerical evaluation of the full Regge action with 20 hinges yields
\begin{equation}
\frac{d^2 S}{d\varphi^2}\bigg|_{\pi/4} \approx -5.9 \times 10^5 < 0,
\end{equation}
confirming that $\varphi = \pi/4$ is a \emph{maximum} of the action
(EXP-047b).  In the Euclidean Regge calculus, the physical
configuration maximizes the action (the Euclidean action is bounded
above, and the dominant contribution comes from the saddle point that
maximizes it).

\medskip\noindent\textbf{Step 4: Conclusion.}
At $\varphi = \pi/4$:
\begin{equation}
|\langle B_1 | B_3 \rangle|^2
= \cos^2\!\left(\frac{\pi}{4}\right)
= \frac{1}{2}. \qedhere
\end{equation}
\end{proof}

\begin{remark}[Analytic vs.\ numerical status]
Steps~1 and~2 constitute a complete analytic proof that
$\varphi = \pi/4$ is stationary.  Step~3 (the sign of
$d^2S/d\varphi^2$) has been verified numerically but not yet proved
analytically.  An analytic proof would require the explicit
$\varphi$-dependence of the Regge action through all 20 hinges.  This
is left as an open problem.
\end{remark}

\begin{remark}[Numerical verification]
EXP-047b scans $\varphi$ over the full range $(0, \pi/2)$ and
confirms that $S(\varphi)$ has a unique maximum at $\varphi = \pi/4$
with the stated second derivative.
\end{remark}

\begin{corollary}[Lattice speed of light]\label{cor:speed_of_light}
Define the lattice speed of light as the inverse of the
$B_1$--$B_3$ overlap squared:
\begin{equation}
c = \frac{1}{|\langle B_1 | B_3 \rangle|^2} = \frac{1}{1/2} = 2 = n_B.
\end{equation}
This recovers the geometric result $c = n_B$ of
eq.~\eqref{eq:c_geometric}, now derived from the variational
principle rather than a density-of-states argument.  The lattice speed
of light equals the number of temporal vertices in each nuclear
simplex.
\end{corollary}

%---------------------------------------------------------------------
\subsection{Implications}
\label{sec:variational_implications}
%---------------------------------------------------------------------

The three theorems supply, from pure geometry, the complete set of
constants entering the closed propagator $P(x) = (1+2x)/(1+x)$:

\begin{center}
\renewcommand{\arraystretch}{1.3}
\begin{tabular}{llll}
\toprule
Theorem & Result & Propagator role & Physical quantity \\
\midrule
1 ($\delta_{AAA} = \pi$) & $e^{i\pi} = -1$ & Alternating sign in $\Sigma$ & Convergent self-energy \\
2 ($\langle\det\rangle = 2/3$) & $\rho = 3/2$ & Impedance per hop & $m_\mu/m_e = \rho/\alpha$ \\
3 ($|\langle B_1|B_3\rangle|^2 = 1/2$) & $c = 2$ & Lattice light speed & Causal structure \\
\bottomrule
\end{tabular}
\end{center}

\noindent
Several features deserve emphasis.

\begin{enumerate}
\item \textbf{Zero free parameters.}
Theorems~1 and~2 depend only on the normalization condition
$\alpha^2 + \beta^2 = 1$; Theorem~3 follows from the $B_1 \leftrightarrow B_2$
exchange symmetry of the Regge action.  No parameter is adjusted.

\item \textbf{Universality.}
All three results are properties of $\partial(\Delta^5)$ with the
$(3,2)$ block structure.  They hold for any realization of the
orthogonality condition $\langle A_i | B_j \rangle = 0$ and any
choice of the A-sector Gram matrix $G_A$.

\item \textbf{Connection to renormalization.}
In conventional quantum field theory, the Dyson equation resums an
infinite series of self-energy insertions.  Theorem~1 shows that the
simplicial geometry \emph{automatically} produces the alternating sign
that makes this series a convergent geometric series.  The closed form
$(1+2x)/(1+x)$ is not an approximation: it is the exact sum, with
convergence guaranteed by $|x| = |\agut \cdot f| < 1$.

\item \textbf{Proof status.}
Theorems~1 and~2 have complete analytic proofs.
Theorem~3 has an analytic proof of stationarity; the maximality
condition ($d^2S/d\varphi^2 < 0$) is verified numerically.
\end{enumerate}

These three results, together with the combinatorial mass formulas of
Chapter~\ref{ch:masses}, yield all charged fermion masses to a median
accuracy of $0.07\%$ and the proton mass to $0.000\%$---with zero
free parameters.
